\documentclass{inprogress}
\newcommand{\tr}{\operatorname{tr}}
\newcommand{\sign}{\operatorname{sign}}
\newcommand{\law}{\operatorname{Law}}
\newcommand{\Span}{\operatorname{Span}}
\newcommand{\erfc}{\operatorname{erfc}}
\newcommand{\eps}{\varepsilon}

\newlist{hypotheses}{enumerate}{2}
\setlist[hypotheses]{leftmargin = 5em}
\newcommand{\Hyp}[1]{\item[\textsc{#1}]\label{hyp:#1}\!\!\!\textsc{:}}
\newcommand{\hyp}[1]{\textsc{#1}}


\title{A compactness approach to $L^2$ stable phase retrieval for independent random variables}

\begin{document}
\maketitle

Let us recall the definition of stable phase retrieval.
\begin{theorem}
    Let $E$ be a subspace of $L^p(\mu)$ and fix $C\geq 1$. We say that $E$ does $C$-SPR if for all $f,g\in E$ we have
    \begin{equation}\label{C-SPR}
        \inf_{|\lambda|=1}\|f-\lambda g\|\leq C \| |f|-|g|\|.
    \end{equation}
\end{theorem}


The purpose of this document is to prove the following theorem.

\begin{theorem}
    \label{thm:goal}
    Let $X_1, \dots, X_n, \dots$ be independent random variables and let $\delta>0$ be such that for all $i\in \mathbb{N}$, we have
    \begin{hypotheses}
        \Hyp{XNorm} $\|X_i\|_2 = 1$,
        \Hyp{XMean} $\mathbb E[X_i] =  0$,
        \Hyp{XEquiv} $\delta \|X_i\|_2 \le \|X_i\|_1 \le (1-\delta) \|X_i\|_2$.
    \end{hypotheses}
    Then the family of $X_i$ performs SPR in $L^2$ with constant $C=C_\delta$. More specifically, for any $X, Y \in \Span(X_1, \dots, X_n \dots)$ one has \[
    \inf_{|\lambda|=1}\mathbb E[|X-\lambda Y|^2]\le C_\delta \mathbb E[||X|-|Y||^2].
    \]
\end{theorem}


We will prove \Cref{thm:goal} by contradiction, together with the standard Taylor-style orthogonality-reduction.


\begin{proposition}\label{by_contradiction}
Let $X_i^n$ be a sequence (of sequences) of random variables, all satisfying \hyp{XNorm}, \hyp{XMean}, \hyp{XEquiv}. Let $a^n$, $b^n$ be sequences in $l^2(\mathbb N)$, and define $X_n:= \sum_i a_i^n X_i^n, Y_n:= \sum_i a_i^n Y_i^n, $. There are no sequences $a^n, b^n$ satisfying all of the following:
\begin{hypotheses}
\Hyp{SNorm} $\|a^n\|_2 = \|b^n\|_2 = 1$, or equivalently $\|X^n\|_{L^2}=\|Y^n\|_{L^2}=1$
\Hyp{SOrth} $\langle a^n, b^n \rangle = 0$, or equivalently $\langle X^n,Y^n\rangle_{L^2}=1$
\Hyp{ASmall} $\| |X^n| - |Y^n| \|_2 \to 0$.
\end{hypotheses}

\end{proposition}

In order to do this, we will use the following compactness/rearrangement result.
\newcommand{\CovFin}{\operatorname{CovFin}}
\begin{lemma}\label{compactness_rearrangement}
    Let $a^n,b^n$ be two sequences of vectors in $\ell^2$ such that $\|a^n\|_2 = \|a^n\|_2 = 1$, $\langle a^n, a^\infty \rangle = 0$. Then, after possible re-indexing of the elements in each $a^n$ (shuffling in the $i$ variable, with a shuffling depending on $n$) and passing to a subsequence in $n$, one can extract sequences  $a^\infty, b^\infty \in L^2$ such that 
    \begin{enumerate}
        \item $a_n \to a^\infty$ in $l^{\infty}$ and $b^n \to b^\infty$ in $l^\infty$
        \item $\lim_{n_0\to\infty} \sup_{n,m\ge n_0} |a_n^m| = 0$ and $\lim_{n_0\to\infty} \sup_{n,m\ge n_0} |b_n^m| = 0$.
    \end{enumerate}
    In this case there exists a symmetric matrix 
    $$\CovFin(a^\infty, b^\infty):= \begin{pmatrix} \langle a^\infty, a^\infty \rangle & \langle a^\infty, b^\infty \rangle \\ \langle b^\infty, a^\infty \rangle & \langle b^\infty, b^\infty \rangle \end{pmatrix}.$$
    By Fatou, $0\preceq \CovFin(a^\infty, b^\infty) \preceq I$.

    Moreover, one can split each sequence $a^n \in l^2$ into a sequence $a^{n} = a^n_{head}+a^n_{tail}$ such that
    \begin{enumerate}
        \item $a^n_{head}$ and $a^n_{tail}$ have disjoint support
        \item $\lim_{n\to\infty} \|a^n_{head}-a^{\infty}\|_{2} = 0$
    \end{enumerate}
    and equivalently for $b_n$. One can do it in such a way that the heads of $a$ and $b$  may have overlapping support, but no head (for $a$ or $b$) overlaps any tail (for $a$ or $b$).
\end{lemma}

\begin{proof}
    For each index $n$, re-order the indices so that $\max (|a_i^n|,|v_i^n|)$ is decreasing. One can then extract a weakly convergent subsequence.
    Then separate \emph{head} from \emph{tail} by a cut-off index $i_0(n)$ that grows slow enough to ensure there is no exess mass.
\end{proof}


The matrix $\CovFin(a^\infty, b^\infty)$ captures how much of the $l^2$-mass escapes to infinity. The part "escaping at infinity" will still be "split" into two parts one coming from an "infinitely divisible" CLT-type term.

From here on, we will assume that the coefficents $a_n$, $b_n$ that we are considering have the property of the Lemma (i.e. that they have been reshufled). This does not change the problem at hand.

Now we will prove some weak facts, 

\begin{lemma} 
    {Phase retrieval}
    Under the hypotnesis of Proposition \ref{by_contradiction}, assuming a rearrangement of the form of Lemma \ref{compactness_rearrangement}, we have
    $
    a^\infty = \pm b^\infty
    $.
\end{lemma}

\begin{proof}
    We can compute $a_i^2$ and $a_i\cdot a_j$ continuously (in $L^1$) from $|X^n|^2$, but $|X^n|^2-|Y^n|^2 \to 0 $ in $L^1$.
\end{proof}

We define the random variables $X^n_{head},X^n_{tail},Y^n_{head},Y^n_{tail}$, noting that by constructions the two heads are independent from the two tails (but the heads are not independent from the heads and vice-versa). The quadruple $(X^n_{head},X^n_{tail},Y^n_{head},Y^n_{tail})$ is tight, and in particular we can extract a convergent subsequence (where all the functions growing strictly slower than quadratically will converge as well). Independence is of course preserved. Moreover,

\begin{lemma}
The tuple $Y^n_{head},Y^n_{tail}$  converges in law to an infinitely-divisible random variable $Y^\infty_{head},Y^\infty_{tail}$. Moreover, the support of $Y^\infty_{head},Y^\infty_{tail}$ is supported on a line if and only if there exists a non-zero vector $(\alpha, \beta)$ such that
$$
\|\alpha a^n - \beta b^n\|_{l^2} \to 0
$$
\end{lemma}

\begin{proof}
    The infinite divisibility follows from the fact that the variables tight and the moments of each of the summands vanish. The variable $(Y^\infty_{head},Y^\infty_{tail})$ has mean $(0,0)$. In particular if it is concentrated on a line, it is a line going through the origin, and there exists a non-zero vector $(s,t)$ such that $\alpha X^\infty_{tail} + \beta+ Y^{\infty}_{tail} = 0$.

    But $\alpha X^\infty_{tail} + \beta+ Y^{\infty}_{tail}$ is equal (in law) to the limit of $Z_n := \sum_{i} (\alpha (a^n_{tail})_i+\beta (b^n_{tail})_i)X_i$. But we have that $0 = \mathbb E [|\lim Z_n|] = \lim \mathbb E [| Z_n|] \approx \lim \mathbb E [| Z_n|^2] $ (TODO this step uses that iid sums preserve up to a constant the $L^/L^2$ ratio)
\end{proof}

By taking limits,

$$
\||X^\infty_{head}+X^\infty_{tail}|-|Y^\infty_{head}+Y^\infty_{tail}|\|_{L^{1}} = 0
$$
this is only possible if the random variable $(X^\infty_{tail}, Y^\infty_{tail})$ is supported on the union of two lines (there must exist some values $x,y$ such that)

$$
|x+X^\infty_{tail}|-|y+Y^\infty_{tail}| = 0 \text{ a. s.}
$$
\begin{lemma}
    If an infinitely-divisible random variable $Z$ in $\mathbb R^2$ is supported on the  union of two secant lines, then it is actually supported in one of the lines. 
\end{lemma}
\begin{proof}
    This already holds for $2-$divisible random variables. If $Z$ is not supported on a line, and $Z  = Y + Y$, then $Y$ cannot be supported on a line. In this case neither can $Y+Y$.
\end{proof}

In particular, in this case we must have $X^{\infty}_{tail} = \pm Y^{\infty}_{tail} \neq 0$. By changing signts, assume equality. Since we also have $X^\infty_{head} = \pm Y^\infty_{head}$, going back to 

$$
|X^\infty_{head}+X^\infty_{tail}| = |\pm X^\infty_{head}+X^\infty_{tail}| \text{ a. s.}
$$
we see that the sign has to be a $+$. 

At this point we have shown that the heads converge in $L^2$ and so do the tails, contradicting orthogonality.



\end{document}